% BatchNorm 平滑损失景观：双子图
% 左：无 BN（高曲率 + 多陡峭 + 尖锐结构）
% 右：有 BN（平滑曲面 + 更宽极小区）
\documentclass[tikz,border=4pt]{standalone}
\usepackage{ctex}
\usepackage{amsmath}
\usetikzlibrary{calc, arrows.meta, decorations.pathmorphing, decorations.pathreplacing, backgrounds}

\begin{document}
\begin{tikzpicture}[scale=1.0, thick]

% ─── 左子图：无 BN（复杂、尖锐）────────────────────────
\begin{scope}[xshift=0cm]
  % 边框
  \draw[black!20, thin, rounded corners=3pt] (-0.3, -0.5) rectangle (6.3, 4.2);
  \node[font=\small\bfseries] at (3.0, 3.9) {无 BatchNorm};

  % 坐标轴
  \draw[->, gray] (-0.1, 0) -- (6.0, 0) node[right, font=\footnotesize] {$\theta$};
  \draw[->, gray] (0, -0.1) -- (0, 3.5) node[above, font=\footnotesize] {$\mathcal{L}$};

  % 凹凸不平、尖锐的损失曲线
  \draw[red!70, very thick]
    plot[smooth, tension=0.6]
    coordinates {
      (0.3, 3.2)
      (0.7, 2.8)
      (1.0, 3.1)
      (1.3, 2.3)
      (1.6, 2.7)
      (2.0, 1.8)
      (2.3, 2.1)
      (2.7, 1.2)
      (3.0, 0.4)
      (3.3, 0.9)
      (3.6, 0.5)
      (3.9, 1.5)
      (4.2, 1.0)
      (4.5, 2.0)
      (4.8, 1.5)
      (5.2, 2.8)
      (5.5, 3.2)
    };

  % 标注尖锐区域
  \draw[<-, red!60, thin] (1.0, 3.1) -- (1.0, 3.4) node[above, font=\tiny, red!60] {尖峰};
  \draw[<-, red!60, thin] (2.3, 2.1) -- (2.3, 2.4) node[above, font=\tiny, red!60] {局部起伏};
  \draw[<-, red!60, thin] (4.2, 1.0) -- (4.2, 1.4) node[above, font=\tiny, red!60] {次极小};

  % 最低点（局部极小）
  \filldraw[red!70] (3.0, 0.4) circle (2.5pt);
  \draw[<-, red!70, thin] (3.0, 0.4) -- (2.5, 0.1) node[below left, font=\footnotesize, red!70] {训练解};

  % 梯度向量（表示高曲率——方向变化快）
  \foreach \x/\y/\dx/\dy in {
    0.8/2.85/0.25/-0.3,
    1.8/2.1/0.2/-0.35,
    2.5/1.55/0.15/-0.4,
    3.5/0.65/0.2/0.3,
    4.0/1.25/0.2/0.25,
    4.7/1.75/0.2/0.4
  }{
    \draw[->, black!40, thin] (\x, \y) -- +(\dx, \dy);
  }
  \node[font=\footnotesize, black!50, rotate=0] at (1.5, 0.3) {梯度方向不稳定};

  % 曲率过大标注
  \draw[decorate, decoration={brace, amplitude=4pt, mirror},
        red!40] (0.6, 0.0) -- (1.4, 0.0);
  \node[below, font=\tiny, red!40] at (1.0, -0.15) {高曲率};

\end{scope}

% ─── 右子图：有 BN（平滑、宽极小）──────────────────────
\begin{scope}[xshift=7.0cm]
  % 边框
  \draw[black!20, thin, rounded corners=3pt] (-0.3, -0.5) rectangle (6.3, 4.2);
  \node[font=\small\bfseries] at (3.0, 3.9) {有 BatchNorm};

  % 坐标轴
  \draw[->, gray] (-0.1, 0) -- (6.0, 0) node[right, font=\footnotesize] {$\theta$};
  \draw[->, gray] (0, -0.1) -- (0, 3.5) node[above, font=\footnotesize] {$\mathcal{L}$};

  % 平滑的损失曲线（BN后）
  \draw[blue, very thick]
    plot[smooth, tension=0.8]
    coordinates {
      (0.3, 3.1)
      (1.0, 2.3)
      (1.8, 1.3)
      (2.6, 0.5)
      (3.2, 0.2)
      (3.8, 0.5)
      (4.6, 1.3)
      (5.4, 2.3)
      (5.7, 3.1)
    };

  % 最低点（宽平极小）
  \filldraw[blue] (3.2, 0.2) circle (2.5pt);
  \draw[<-, blue, thin] (3.2, 0.2) -- (3.2, -0.15) node[below, font=\footnotesize, blue] {训练解};

  % 宽平极小区域标注
  \draw[blue!40, dashed, thin] (2.0, 0) -- (2.0, 0.45);
  \draw[blue!40, dashed, thin] (4.4, 0) -- (4.4, 0.45);
  \draw[blue!40, <->, thin] (2.0, 0.35) -- (4.4, 0.35)
    node[midway, above, font=\footnotesize, blue!70] {宽极小区域};

  % 梯度向量（方向一致、平滑）
  \foreach \x/\y/\dx/\dy in {
    0.7/2.95/0.25/-0.35,
    1.4/1.95/0.25/-0.4,
    2.2/0.95/0.2/-0.3,
    4.2/0.95/0.2/0.3,
    5.0/1.95/0.2/0.35,
    5.5/2.6/0.2/0.4
  }{
    \draw[->, black!40, thin] (\x, \y) -- +(\dx, \dy);
  }
  \node[font=\footnotesize, black!50] at (1.5, 0.25) {梯度方向稳定};

  % 曲率平缓标注
  \draw[decorate, decoration={brace, amplitude=4pt, mirror},
        blue!40] (2.0, 0.0) -- (4.4, 0.0);
  \node[below, font=\tiny, blue!40] at (3.2, -0.15) {低曲率};

  % Lipschitz 梯度标注
  \node[font=\footnotesize, blue!70, align=center] at (3.2, 3.2)
    {$\|\nabla_\mathbf{w}\mathcal{L}_\text{BN}\|$\\$\leq \|\nabla_\mathbf{w}\mathcal{L}_\text{no-BN}\|$};
\end{scope}

% 中间对比箭头
\draw[->, black!60, very thick] (6.35, 1.8) -- (7.0, 1.8)
  node[midway, above, font=\footnotesize, black!60] {};
\node[font=\small, black!60] at (6.65, 2.1) {加 BN};

\end{tikzpicture}
\end{document}
